\section{Artificial Intelligence}
\label{sec:artificial-intelligence}

AHI combines three kinds of artificial intelligence, each with a different job: statistical models that establish what is normal for each person, machine learning models that detect patterns and estimate risk, and a large language model (LLM) that writes the results in plain, direct language. The LLM is accessed through the \textbf{Claude API} from Anthropic \cite{anthropic_claude_api} whenever there is a link to Earth, and through a \textbf{local model specialized in AHI}, running on board, when there is not. The central rule is that \textbf{generative AI explains the results but is never responsible for the clinical calculations}. Status, risk, forecasts and the choice of recommendations come from the analytical engine of \ref{sec:analytical-methodology}, where every output can be traced to data and to a versioned model or rule.

\subsection{Roles of Each Layer}
\label{subsec:ai-roles}

\begin{table}[htbp]
\centering
\vspace{0.8em}
\caption{The layers of artificial intelligence in AHI.}
\label{tab:ai-layers}
\small
\begin{tabular}{p{3.4cm}p{5.0cm}p{5.6cm}}
\hline
\textbf{Layer} & \textbf{Methods} & \textbf{Responsibility} \\
\hline
Statistical & Robust baselines, drift detection, damped trend projection, explicit rules & Personal normal, acute and slow deviations, first forecasts; works from the first day \\
Machine learning & Isolation Forest, XGBoost, random forests; later LSTM and Temporal Fusion Transformer & Multivariate anomalies, risk estimates, factor ranking, multivariate forecasts \\
Language model, connected & Claude API with retrieval of curated sources & Plain language explanations of predictions, alerts and recommendations; summaries and handover text \\
Language model, on board & A local model specialized in the AHI context & The same explanations when the link to Earth is lost \\
\hline
\end{tabular}
\vspace{1.2em}
\end{table}

\subsection{Machine Learning Models}
\label{subsec:ml-models}

\medskip\noindent\textbf{Isolation Forest} isolates observations that are unusual in several signals at once, which complements the univariate checks. It is trained per person or per small group, on data that the baselines already classify as normal.

\medskip\noindent\textbf{XGBoost and random forests} estimate the risk of a defined event within a window, such as a fatigue episode in the next 72 hours, from features such as sleep debt, heart rate variability relative to baseline, training load and environmental exposure. Their feature importance and per prediction contributions identify which factors drove a result.

\medskip\noindent\textbf{Sequence models}, LSTM networks and the Temporal Fusion Transformer, forecast several signals together. They need long histories and are added only when such histories exist.

\medskip\noindent\textbf{The data problem.} Learned models need data, and data on astronauts are scarce: there is no public continuous telemetry, crews are small, and labeled outcomes are rare. AHI therefore follows a staged approach. It starts with baselines and documented rules, which need no training. It adds Isolation Forest and XGBoost when enough data from the mission and from authorized historical sources exist, and sequence models later. Public NASA data, analog studies and simulated data are used to develop and test the models, never as a substitute for validation on the target population (\ref{sec:data-sources}). A model trained only on simulated data learns the simulator and is not used for crew decisions.

\medskip\noindent\textbf{Governance.} Each model has a model card with its purpose, features, training data, version, evaluation results and known limits. Each output records the model version that produced it. A new model version runs in shadow mode, next to the current one and without affecting the crew view, until its evaluation is reviewed.

\subsection{Claude API Explanation Layer}
\label{subsec:llm-layer}

Statistical predictions are precise but hard to read under pressure. A line such as ``projected resting heart rate 71 beats per minute at 48 hours, 80\% interval 66 to 76, band exit in 31 hours, confidence 0.72'' is correct, but it does not tell a crew member what it means. AHI uses the Claude API to rewrite such outputs in a more direct way, for example: ``If the current trend continues, your resting heart rate is likely to rise above your usual range in about a day. The estimate is moderately certain because two nights of sleep data are missing.'' The numbers in the sentence are the numbers of the engine, and nothing is added.

The tasks of the layer are the following:

\begin{itemize}
\item \textbf{Write predictions directly:} turn each forecast, with its horizon, interval and confidence, into one or two plain sentences.
\item \textbf{Explain alerts and recommendations} already selected by the engine: what changed, compared with what, and why it matters.
\item \textbf{Summarize} the condition of a crew member or of the crew for a shift handover.
\item \textbf{Answer questions} about the data of the platform, citing the records and the curated scientific sources it used.
\item \textbf{Adapt to the reader:} plain language for the crew, technical detail for the flight surgeon.
\end{itemize}

The API receives only structured outputs of the engine, such as the metric, the value, the baseline, the forecast, the factors, the confidence and the recommendation, with pseudonymized identifiers and passages from a curated library of NASA and peer reviewed sources. It does not receive raw signals, names or identifying data, and it does not compute status or risk. The model in use is recorded with every generated text. The choice of model follows the current Claude models available through the API, and it is revised as new models are released.

\subsection{Operation Without Contact with Earth}
\label{subsec:offline-llm}

The Claude API is an internet service. It is reachable only while the spacecraft has a link to Earth. Space links are interrupted: during Artemis~II, the Orion spacecraft lost all contact with Earth for about 40 minutes while it passed behind the Moon \cite{ctv_artemis2_blackout}, and the communication delay grows with distance up to about 24 minutes one way for Mars \cite{nasa_space_communications_seven_things}. Interruptions are also possible on the ISS. Health monitoring cannot stop during these periods, and it is exactly when the crew is alone that a clear explanation matters most.

The analytical engine already runs on board and does not depend on the link (\ref{sec:software-architecture}). For the explanation layer, the proposed solution is to \textbf{train a local language model specialized in the context of AHI}, installed on the computers of the ISS or of the Artemis spacecraft and available to the astronauts without any connection. Its design follows four principles:

\begin{itemize}
\item \textbf{Specialized and small.} The local model does not need general knowledge. It needs to explain the outputs of AHI: the metric catalog, the eight health dimensions, the rules, the forecasts, the recommendation catalog and the curated scientific library. An open model small enough for the on board computers is adapted to this context with examples written and reviewed by the team and by flight surgeons.
\item \textbf{Same inputs, same guardrails.} It receives the same structured outputs as the Claude API and its texts go through the same checks (\ref{subsec:guardrails}).
\item \textbf{Automatic switch.} When the link is lost, the platform switches to the local model and labels its texts as generated on board. When the link returns, texts can again be written through the Claude API, and the history is synchronized with the ground.
\item \textbf{Validated before flight.} The local model is evaluated on reference cases before it is installed, for factual consistency, absence of diagnostic wording and correct use of citations, and a new version is installed only after review.
\end{itemize}

If neither model is available, the platform shows the structured result with template text. Monitoring, alerts and recommendations never depend on a language model.

\subsection{Guardrails}
\label{subsec:guardrails}

The same rules apply to the Claude API and to the local model:

\begin{itemize}
\item \textbf{No diagnosis and no prescription.} The model does not name diseases, does not suggest medication or doses, and does not change the level or the recommendation selected by the engine.
\item \textbf{Grounded statements only.} Every number in a generated text must exist in the structured input, and every scientific statement must cite a passage from the curated library. A text that fails these checks is not shown, and the structured result is shown instead.
\item \textbf{No invented limits.} The model never states a threshold or an official limit that is not in its input.
\item \textbf{Uncertainty kept visible.} Confidence, missing data and limitations from the input must appear in the text, and a forecast is never written as a certainty.
\item \textbf{Origin visible.} Generated text is labeled as generated, with the model and the place where it ran, and the structured result is always one click away.
\end{itemize}

\subsection{Open Decisions}
\label{subsec:ai-open-decisions}

Four decisions remain open and need input from mission medical and security staff. The first is the data agreement for the Claude API: which structured, pseudonymized data may leave the spacecraft, and under which retention terms. The second is the choice of the open model for the local layer and of the on board hardware that runs it. The third is how the local model is updated during a long mission, given limited bandwidth. The fourth is the shared evaluation of both language layers, with a set of reference cases that checks factual consistency, the absence of diagnostic wording and the correct use of citations.
